% !PRINTED 545
\weberpage{545}
We state it thus:

1. \gesperrt{If $\pi$ is a prime divisor of the natural prime $p$, then for every integer $\omega$ in the field $\Omega$}
\[ \omega^{N_a(\pi)} \equiv \omega \pmod{\pi}. \]

Important consequences now attach to this:

2. \gesperrt{If $f(t)$ denotes an integral function of degree $m$ whose coefficients are integers in $\Omega$, and $\pi$ a prime functional, then the congruence}
\weq{5}{f(t) \equiv 0 \pmod{\pi}}
\gesperrt{has at most $m$ roots, that is, there are at most $m$ incongruent integers in $\Omega$ which, put for $t$, satisfy the congruence.}

For if $\alpha$ denotes any number in $\mathfrak{o}$, we can set
\weq{6}{f(t) = (t-\alpha)\,f_1(t) + f(\alpha)}
in which $f_1(t)$ is just such a function as $f(t)$, but only of degree $(m-1)$. But if $f(\alpha) \equiv 0 \pmod{\pi}$, then every root of (5) must satisfy the congruence
\[ (t-\alpha)\,f_1(t) \equiv 0 \pmod{\pi} \]
It must therefore either be congruent with $\alpha$, or be a root of the congruence of degree $(m-1)$
\[ f_1(t) \equiv 0 \pmod{\pi} \]
If we assume our theorem as proved for congruences of degree $(m-1)$, then it accordingly holds also for congruences of degree $m$; and since it holds for congruences of the first degree, it is generally true.

Every number in $\mathfrak{o}$ not divisible by $\pi$ satisfies, as we have seen, the congruence
\weq{7}{\omega^{p^f-1} \equiv 1 \pmod{\pi}.}
Now if $a$ is the least natural number for which the congruence
\weq{8}{\omega^a \equiv 1 \pmod{\pi}}
is satisfied, then it can be shown by the method of inference already often applied that every other exponent $l$ for which $\omega^l \equiv 1$ must be a multiple of $a$. For if $l$ were not divisible by $a$, then, if $a'$ is the residue of the division of $a$ by $l$, we should also have $\omega^{a'} \equiv 1$, which by the assumption about $a$ is not possible. % [sic] the source has ``division of $a$ by $l$''; the minimality argument requires the remainder of $l$ by $a$
Hence $a$ is a divisor of $p^f-1$, and we call $\omega$ a \gesperrt{number belonging to the exponent $a$}.
